\documentclass[mode = thesis, fontset = lxgw]{westlakethesis}

\usepackage{hologo}

\westlakeset{
  id        = 20250000000,
  school    = sls,
  title     = 西湖大学学位论文模板 (基于 \hologo{LaTeX3} 开发)/
              Westlake University Thesis Template Constructed with
              \hologo{LaTeX3},
  subject   = 生物学/Biology,
  author    = 鼠小西/West Mouse,
  PI        = 果小蝇/Fruit Fly,
  % master  = false,
  bibsource = reference,
}

\begin{document}

\maketitle

\frontmatter
\begin{abstract}[cn]
  随着 XXXXX 的发展，XXX 问题日益突出。本文提出了一种新型算法模型，融合了三项关键技术，
  通过实验设计，在数据集上验证了方法的有效性，取得了指标的性能提升。研究表明，
  该方法能够明显优势，为该问题提供了新的解决方法。
  XXX XXX XXX XXX XXX XXX XXX XXX XXX XXX XXX XXX XXX XXX XXX
  XXX XXX XXX XXX XXX XXX XXX XXX XXX XXX XXX XXX XXX XXX XXX
  XXX XXX XXX XXX XXX XXX XXX XXX XXX XXX XXX XXX XXX XXX XX.\par
  \zhlipsum[1]
  \keywords{关键词 1, 关键词 2, 关键词 3, 关键词 4, 关键词 5（不超过 5 个）}
\end{abstract}

\begin{abstract}[en]
  With the rapid development of financial technology (FinTech), traditional
  credit risk assessment models face challenges in processing large-scale,
  heterogeneous data. This study proposes a novel ensemble learning framework
  that combines deep neural networks with interpretable machine learning models
  to improve prediction accuracy while maintaining transparency.
  XXX XXX XXX XXX XXX XXX XXX XXX XXX XXX XXX XXX XXX XXX XXX
  XXX XXX XXX XXX XXX XXX XXX XXX XXX XXX XXX XXX XXX XXX XXX
  XXX XXX XXX XXX XXX XXX XXX XXX XXX XXX XXX XXX XXX XXX XX.\par
  \lipsum[1]
  \keywords{Keywords 1, Keywords 2, Keywords 3}
\end{abstract}

\tableofcontents
\mainmatter

\chapter{章标题}

学位论文正文内容（小四号字，宋体，Times New Roman 体，两端对齐书写，
段落首行左缩进 2 个汉字符）\cite{yang2009cernet}。

\section{一级节标题}

学位论文正文内容（小四号字，宋体，Times New Roman 体，两端对齐书写，
段落首行左缩进 2 个汉字符）。

XXX XXX XXX XXX XXX XXX XXX XXX XXX XXX XXX XXX XXX XXX XXX
XXX XXX XXX XXX XX\cite{bonjour2009minerals}.

\subsection{二级节标题}

学位论文正文内容。

XXX XXX XXX XXX XXX XXX XXX XXX XXX XXX XXX XXX XXX XXX XXX
XXX XXX XXX XXX XX.

\subsubsection{三级标题}

学位论文正文内容。

XXX XXX XXX XXX XXX XXX XXX XXX XXX XXX XXX XXX XXX XXX XXX
XXX XXX XXX XXX XX.

\section{本文研究主要内容}

\section{本文研究意义}

\section{本章小结}

\chapter{章标题}

\section{论文正文}

学位论文正文内容。

XXX XXX XXX XXX XXX XXX XXX XXX XXX XXX XXX XXX XXX XXX XXX
XXX XXX XXX XXX XX.

\section{字数要求}

\subsection{硕士论文要求}

\subsection{博士论文要求}

\section{其他要求}

\subsection{页面设置}

\section{本章小结}

\chapter{章标题}

\section{图表格式}

\section{公式格式}

\section{本章小结}

\chapter{章标题}

\section{主要结论}

\section{研究展望}

\printbibliography

\chapter*{致谢}

\commitment[
  example-image-a.pdf, 2026 年 5 月 20 日,
  example-image-a.pdf, 2026 年 5 月 21 日,
  example-image-b.pdf, 2026 年 5 月 21 日,
]

\end{document}